%% notes-tex/031-unified-representation/sections/10-representation.tex

\section{The representation\secstatus{todo}}
\label{sec:representation}

Figure~\ref{fig:unified} is the input representation and the decoder the preceding sections
argue for, drawn as the tensors each channel is. It has two per-record input channels and a
per-record token, one shared molecule encoder, one operator and one decoder, and it carries
its own limits.

%% SOURCE: notes/experiments.031-env-chemgen-inhibitor-tolerance.mermaid.unified-input.md,
%% rendered by notes/assets/publish/scripts/mermaid_pdf.sh through `make diagrams`.
\tcfig{figures/unified_input.pdf}{One input representation and one decoder for the four kept
datasets, in the notation of the cell graph transformer: $N$ genes, $d_h$ the hidden width,
$S$ the perturbed set, $H_{\mathrm{pert}}$ the Type I output, $R_\phi$ a Type II readout.
Gray is a served dataset, purple a per-record input channel, yellow the fixed embedding table,
orange a computation, blue the optional metabolism branch, red what the representation drops
or does not carry. The genotype is a per-gene functional dose whose background level is the
ploidy, and it scales the perturbation operator through $\gamma(m_j)$. The environment is a
compound set and a medium set through the shared molecule encoder, a dose-basis token and a
masked log molar value, pooled into one context $z_E$. The dataset token $t$ carries the
measurement scale and every physical field, because each of those is constant within a
dataset. The decoder is one scalar head over the concatenation.}{fig:unified}

\notesec{Genotype}
A per-gene functional dose over the gene universe, $d \in \{0, \tfrac{1}{2}, 1\}^{N}$ with
$d_i = c_i / r_i$, copies present over copies in the reference. Zero for a deletion, one half
for a heterozygous deletion in a diploid, one for an unperturbed gene. Ploidy needs no
separate feature because it is the background level. The perturbed set is $S = \{i : d_i <
1\}$ and the magnitude $m_i = 1 - d_i$ is what scales the operator. This is lossless for these
datasets and is validated by Hoepfner's 648{,}977 paired cells.

\notesec{Environment}
The set of dosed compounds as InChIKeys through the shared molecule encoder $f$, the medium
components through the same encoder, a dose-basis one-hot over target inhibition, fixed
concentration and unstated, and the log molar concentration where a record has one with a
mask where it does not. The basis token is the honest encoding given Section~\ref{sec:dose}:
a target-inhibition dose, a fixed molar dose and an unstated dose are three different
statements, and a single number would assert a comparability that does not exist. The
masked log molar value keeps the within-dataset dose variation that Hillenmeyer and Hoepfner
do carry without pretending it compares across them. The four parts are concatenated into one
context $z_E$, which enters at the readout in the first arm and as extra key-value rows of
the operator in the second (Section~\ref{sec:model}).

\notesec{Dataset token}
A one-hot over the four sources, projected and concatenated at the readout. It carries the
response scale, the assay, and every physical field, since temperature, aerobicity, duration,
pH and solvent are constant inside each dataset (Table~\ref{tab:physical}). No separate
physical channel exists for that reason; the fields stay in the record.

\notesec{The encoder table}
One embedding table, keyed by InChIKey, serves the dosed compounds, the medium components and
the Yeast9 metabolites. It is precomputed by \file{embed\_compounds.py},
\file{embed\_media\_components.py} and \file{yeast9\_molecule\_coverage.py} for all twelve
encoders, so the encoder is an arm selected at data-loading time and never a training-time
cost. FCFP4 count is the default because it is the measured best on the compound-cold target
(Table~\ref{tab:baselines}); the transformers and MolE are the alternatives that a first
round compares against it.

\notesec{What it does not carry}
The dose vector drops cassette identity, barcode and collection. Dose is never one number
across datasets. Alleles off the dosage axis are out of scope for this phase. And 167 Yeast9
species have no structure and need a non-molecular slot. Each of those is a measured limit
rather than an anticipated one, and each bounds a specific later step: the dose limit bounds
any cross-dataset dose-response claim, and the structureless species bound the metabolic
branch.

\notesec{Implementation, in order}
\begin{enumerate}
\item The dataset-to-tensor processor emits, beside the perturbed gene indices, a per-index
  magnitude $m_j$ computed from the perturbation's copy number and reference copy number, and
  the reference genome's ploidy; a deletion emits one. Every existing dataset produces the
  same tensors as today plus a vector of ones.
\item The perturbation operator takes the magnitude and scales each key-value row by
  $\gamma(m_j)$ with $\gamma(1) = 1$, first as the fixed identity $\gamma(m) = m$ and then as
  a learned scalar map; the flux layer reads $d_i$ in place of its presence bit.
\item The processor emits the environment context from the served record: the compound
  InChIKeys and medium component keys resolved against the precomputed embedding table, the
  dose basis as a one-hot, and the log molar value with its mask. A record whose compound has
  no vector in the table raises rather than receiving zeros.
\item The processor emits the dataset index over a fixed vocabulary of the four sources, and
  the label transform standardizes each source's response within the training split.
\item The head takes the concatenation of the CLS embedding, the pooled perturbed-gene
  summary, the environment context and the projected dataset token, as one scalar head under
  the masked loss.
\item The evaluation is fixed before the first run: compound-cold within Vanacloig, 41 folds,
  scored on the centered and the raw target with the reliability ceiling beside every
  compound, and a Vanacloig-only training arm as the control for pooling.
\end{enumerate}

\notesec{First round}
Four arms on the fixed evaluation, three seeds each: Vanacloig alone; the four datasets with
$z_E$ at the readout; the four datasets with $z_E$ in the operator; and the operator arm with
the learned $\gamma$. The first two answer whether pooling helps at all, the third whether the
gene-by-compound interaction needs the operator, and the fourth whether heterozygosity is an
attenuation. The graph prior runs at $\lambda = 1$ throughout, and the annealed mask is a
fifth arm if the round has budget for it. Every score is reported against the ceiling of its
compound, and a partial run is a partial run.
